% Source pp. 254-263 (Lectures19-25.pdf).
\section{Operator Formalism vs Functional Integral in QFT}

So far we have discussed QFT mostly from the perspective of (Euclidean space-time)
functional integral. There are aspects which are better understood within the
formalism based on the notions of the Hilbert space of states which supports the
algebra of local field operators. In the previous discussion I made occasional
references to statements that follow from the operator formalism, mostly from the
positivity of norms in the Hilbert space, as well as from the so called spectral
condition. Here I would like to recollect the basics of that formalism, and
clarify its relation to the functional integral.

\subsection*{Operator formalism}

As was stated before (Section~7), the correlation functions defined by the
functional integrals
\begin{equation}\label{eq:corrPhi}
\bigl\langle O_1(x_1)\cdots O_N(x_N)\bigr\rangle:=Z^{-1}\int[D\Phi]\,O_1(x_1)\cdots O_N(x_N)\,e^{-\mathcal A[\Phi]}
\end{equation}
can be interpreted as the vacuum-vacuum matrix element of the ``T-ordered'' product
\begin{equation}\label{eq:Tordered}
\bigl\langle O_1(x_1)\cdots O_N(x_N)\bigr\rangle=\bra0T_\tau\bigl(\hat O_1(x_1)\cdots\hat O_N(x_N)\bigr)\ket0
\end{equation}
of the Heisenberg field operators $\hat O(x)$. The latter notions involve
nominating one of the coordinates to the role of the ``Euclidean time'' $\tau$, so
that $x=(\xx,\tau)$. Usually, the Heisenberg picture is developed in the real time
$t$ which relates to $\tau$ by the analytic continuation to $\tau=it$, but in most
of the development below this specification is not necessary; we generally will
regard $\tau$ as a complex number. Then
\begin{equation}\label{eq:Heis}
\hat O(\xx,\tau)=e^{\hat H\tau}\,\hat O(\xx)\,e^{-\hat H\tau}\,.
\end{equation}
By construction, the field operators satisfy Heisenberg ``equation of motion''
\begin{equation}\label{eq:HeisEOM}
\frac\partial{\partial\tau}\,\hat O(\xx,\tau)=[\hat H,\hat O(\xx,\tau)]
\end{equation}
In fact, most of the calculations can be done at real $\tau$ (``imaginary time'').
In that case the ``time ordering'' $T_\tau$ means that the operators in
\eqref{eq:Tordered} are ordered according to decreasing values of $\tau$: the
operators with greater $\tau$ are placed to the left.

We have derived this relation in quantum mechanics (Section~5), and then assumed
that it can be extended to QFT. In fact, such extension is not straightforward,
and involves many subtleties mostly related to regularization and renormalization
(our brief inspection of the Wilson's RG gives the idea of the size of the
problem). Here I would like to revisit the relation \eqref{eq:Tordered} from more
general perspective. Starting from the functional integral definition of the
correlation functions, we will try to recover basic notions of the operator
formalism. Eq.~\eqref{eq:Heis} involves the Hamiltonian, and we would like to see
how the Hamiltonian, or more generally, the Energy-Momentum $P^\mu$, can be
introduced in terms of the correlation functions.

\subsection*{Correlation Functions and Energy-Momentum}

Consider the correlation function
\[
\bigl\langle T^{\mu\nu}(x)\,O_1(x_1)\cdots O_N(x_N)\bigr\rangle
\]
involving the energy-momentum tensor field $T^{\mu\nu}(x)$. As we have seen in the
previous Section, it satisfies the Ward Identity \eqref{eq:Ward2}, in particular
\[
\bigl\langle\partial_\mu T^{\mu\nu}(x)\,O_1(x_1)\cdots O_N(x_N)\bigr\rangle=0\qquad\text{for}\quad x\neq x_1,\dots,x_N\,,
\]
As a consequence, the integral
\[
\int_\Sigma d\Sigma_\mu(x)\,\bigl\langle T^{\mu\nu}(x)\,O_1(x_1)\cdots O_N(x_N)\bigr\rangle
\]
taken over a $(d-1)$-hypersurface $\Sigma$ does not change under topologically
trivial deformations of $\Sigma$. The hypersurface $\Sigma$ divides the space
$\R^d$ into two parts, and correspondingly splits the set of insertion points
$x_1,\dots,x_N$ into two subsets, according to their locations within the parts.
Topologically trivial are those deformations which do not change this splitting,
or, in other words the deformations which do not make the hypersurface ``cross
over'' any of the insertion points. For the sake of our present discussion it is
convenient to choose $\Sigma$ to be an ``equal-time slice''. As before, we
nominate one of the coordinates in $\R^d$ to be the (imaginary) ``time'' $\tau$,
i.e.\ $x=(\xx,\tau)$, and define $\Sigma(\tau)$ as the union of points with given
$\tau$. Define\footnote{The factor $i$ is added to this notation to agree with the
conventional notation $P^\mu$ for the energy-momentum in real Minkowski space time.
Upon continuation $\tau=it$ the components of the energy-momentum tensor transform
so that $T^{\tau k}=iT^{0k}$, $k=1,2,\dots,d-1$, and $T^{\tau\tau}=-T^{00}$.}
\begin{equation}\label{eq:Pslice}
iP^\mu(\tau):=\int_{\Sigma(\tau)}d\Sigma_\nu(x)\,T^{\mu\nu}(x)=\int_{\Sigma(\tau)}d^{d-1}\xx\,T^{\mu\tau}(\xx,\tau)\,,
\end{equation}
and assume that
\begin{equation}\label{eq:tauorder}
\tau_1>\tau_2>\dots>\tau_M>\tau>\tau_{M+1}>\dots>\tau_N\,,
\end{equation}
as illustrated in figure~\ref{fig:slice}.
\begin{figure}[ht]
\centering\figfile{slice}
\caption{The equal-time slice $\Sigma(\tau)$ separates the insertion points with
$\tau_i>\tau$ from those with $\tau_i<\tau$.}
\label{fig:slice}
\end{figure}
Then the correlation function
\begin{equation}\label{eq:Pcorr}
\bigl\langle P^\mu(\tau)\,O_1(x_1)\cdots O_N(x_N)\bigr\rangle
\end{equation}
does not change if one changes the value of $\tau$, as long as the order
\eqref{eq:tauorder} is preserved. We therefore drop the argument $\tau$ writing in
the notation for $P^\mu:=P^\mu(\tau)$, but adopt the convention that the position of
the insertion $P^\mu$ in the correlation function corresponds to the order
\eqref{eq:tauorder}, i.e.\ \eqref{eq:Pcorr} is written as
\begin{equation}\label{eq:Pinsert}
\bigl\langle O_1(x_1)\cdots O_M(x_M)\,P^\mu\,O_{M+1}(x_{M+1})\cdots O_N(x_N)\bigr\rangle
\end{equation}
With this convention the position of the insertion $P^\mu$ represents the
topological class of the surface $\Sigma$ in \eqref{eq:Pslice} relative to the
points $x_i$.

Now that the position of $P^\mu$ in the correlation function \eqref{eq:Pcorr}
matters, one may ask about the ``commutators'' of $P^\mu$ with the local field
insertions. We have formally (see figure~\ref{fig:commslab})
\begin{multline*}
iP^\mu O(\xx,\tau)-O(\xx,\tau)\,iP^\mu\\
=\int d^{d-1}\yy\,T^{\tau\mu}(\yy,\tau+\Delta\tau)\,O(\xx,\tau)-\int d^{d-1}\yy\,T^{\tau\mu}(\yy,\tau-\Delta\tau)\,O(\xx,\tau)\,,
\end{multline*}
\begin{figure}[ht]
\centering\figfile{commslab}
\caption{The commutator with $P^\mu$ is the difference of the slices at
$\tau+\Delta\tau$ and $\tau-\Delta\tau$, which can be closed into a small surface
around $x$.}
\label{fig:commslab}
\end{figure}
or, by closing the hypersurface
\begin{equation}\label{eq:Pcomm}
i[P^\mu,O(x)]=\int_{\Sigma_x}d\Sigma_\nu(y)\,T^{\mu\nu}(y)\,O(x)
\end{equation}
where $\Sigma_x$ is an arbitrary small closed hypersurface wrapped around the point
$x$. Again, this relation is understood as being valid for the insertions in the
correlation functions, with the above ordering convention.

The integral in the r.h.s.\ can be calculated as follows. Consider again the Ward
identity \eqref{eq:Ward},
\begin{equation}\label{eq:Ward3}
\sum_{i=1}^n\bigl\langle\mathcal O_1(x_1)\cdots\delta_\varepsilon\mathcal O_i(x_i)\cdots\mathcal O_n(x_n)\bigr\rangle
=-\int d^dx\,\partial_\mu\varepsilon_\nu(x)\,\bigl\langle T^{\mu\nu}(x)\,\mathcal O_1(x_1)\cdots\mathcal O_n(x_n)\bigr\rangle\,.
\end{equation}
where $\delta_\varepsilon O(x)$ stand for the variations of the fields $O(x)$ under
the coordinate transformation $x^\mu\to x^\mu+\varepsilon^\mu(x)$. The integration
in the r.h.s.\ is over the whole space $D:=\R^d$. Let us split this domain into
several pieces, as follows
\[
D:=\R^d=\sum_{i=1}^ND_i+D_{\text{out}}\,,
\]
where $D_i$ is a small domain containing the point $x_i$, and $D_{\text{out}}$ is
the remaining part of $D$ left after removing all $D_i$\footnote{In conventional
set theory symbolics the decomposition reads
$D_{\text{out}}:=D\setminus(\cup_iD_i)$.} (figure~\ref{fig:Di}).
\begin{figure}[ht]
\centering\figfile{Di}
\caption{Left: a small domain $D_i$ around the insertion point $x_i$. Right: the
normal element $d\Sigma_\mu$ of its boundary $\partial D_i$.}
\label{fig:Di}
\end{figure}
The integral splits accordingly, as
\[
\int d^dx=\sum_{i=1}^N\int_{D_i}d^dx+\int_{D_{\text{out}}}d^dx
\]
For all $x\in D_{\text{out}}$ we have
\[
\partial_\mu T^{\mu\nu}(x)=0\qquad\text{for}\quad x\in D_{\text{out}}\,,
\]
and hence the integral over $D_{\text{out}}$ can be transformed via the Stokes
theorem, as
\[
\int_{D_{\text{out}}}d^dx\,\partial_\mu\bigl\langle\bigl(\varepsilon_\nu(x)\,T^{\mu\nu}(x)\bigr)\,O_1(x_1)\cdots O_N(x_N)\bigr\rangle=-\sum_{i=1}^NI_i\,,
\]
where $I_i$ stand for the surface integrals over the boundaries of $D_i$,
\[
I_i=\int_{\partial D_i}d\Sigma_\mu(x)\,\varepsilon_\nu(x)\,\bigl\langle T^{\mu\nu}(x)\cdots O_i(x_i)\cdots\bigr\rangle
\]
and $d\Sigma_\mu(x)$ is the normal element of the boundary $\partial D_i$. Hence,
the r.h.s.\ of the Eq.~\eqref{eq:Ward3} splits into the sum
\begin{equation}\label{eq:splitsum}
\sum_{i=1}^N\Bigl(I_i-\int_{D_i}d^dx\,\partial_\mu\varepsilon_\nu(x)\,\bigl\langle T^{\mu\nu}(x)\cdots\bigr\rangle\Bigr)\,.
\end{equation}
By construction, each of the terms here is independent of the shape of the domain
$D_i$, which can be made arbitrary small. This implies that each of these terms
depends on $\varepsilon(x)$ in a local way, i.e.\ the term associated with $D_i$
depends on $\varepsilon(x_i)$, $\partial\varepsilon(x_i)$, $\dots$. Comparing
\eqref{eq:splitsum} with the l.h.s.\ of \eqref{eq:Ward3} we conclude that the
equality is satisfied term by term in the sum over $i$. This leads to the following
representation of the variations $\delta_\varepsilon O(x)$
\begin{equation}\label{eq:deltaO}
\delta_\varepsilon O(x)=\int_{\partial D_x}d\Sigma_\mu(y)\,\varepsilon_\nu(y)\,T^{\mu\nu}(y)\,O(x)-\int_{D_x}d^dy\,\partial_\mu\varepsilon_\nu(y)\,T^{\mu\nu}(y)\,O(x)\,.
\end{equation}
The relation is understood as valid for the correlation functions. Let me note again
that the r.h.s.\ is in fact independent of the choice of the (small) domain $D_x$.

As was already mentioned, the variation $\delta_\varepsilon O(x)$ admits expansion in
the derivatives of $\varepsilon$ taken at the point $x$, with the leading term
\[
\delta_\varepsilon O(x)=\varepsilon^\mu(x)\,\partial_\mu O(x)+\text{terms involving }\partial\varepsilon(x),\dots
\]
In particular, for a constant $\varepsilon^\mu(x)=\varepsilon^\mu$ the terms with
the derivatives do not contribute. On the other hand, in this case the second term
in \eqref{eq:deltaO} vanishes, and we obtain
\[
\partial^\mu O(x)=\int_{\partial D_x}d\Sigma_\nu(y)\,T^{\mu\nu}(y)\,O(x)
\]
The r.h.s.\ here coincides with the expression in \eqref{eq:Pcomm}. Finally, we
find
\[
[P^\mu,O(x)]=-i\,\partial^\mu O(x)\,.
\]
Taking $\mu$ to be the ``time'' index $\tau$, and recalling that $P^\tau=iP^0=iH$,
one recovers \eqref{eq:HeisEOM}.\footnote{Up to the overall sign. With the sign
conventions used in this and the previous Section one finds
$\partial_\tau O=-[H,O]$; the sign in \eqref{eq:HeisEOM} is recovered by reversing
the overall sign of $T^{\mu\nu}$ in \eqref{eq:Tdef}. The notes flag this sign
without resolving it.} The ``spatial'' components of this equation
\[
i\nabla O(\xx,\tau)=[\PPv,O(\xx,\tau)]
\]
reproduce the well known (see e.g.\ Chapter 2 of PS) representation of the
Heisenberg field operators
\[
\hat O(\xx,t)=e^{-i\PPv\cdot\xx+iHt}\,\hat O(\mathbf0,0)\,e^{i\PPv\cdot\xx-iHt}
\]
which I chose here to write directly in Minkowski space-time.

\subsection*{Space of States\protect\footnote{The ideas displayed in this section
are borrowed from the construction in K.~Osterwalder, R.~Schrader, ``Axioms for
Euclidean Greens functions,'' Comm.\ Math.\ Phys.\ \textbf{31} (1973), 83--112;
\textbf{42} (1975), 281--305.}}

The functional integral is particularly suitable for dealing with the correlation
functions \eqref{eq:corrPhi}. Calculations in the previous paragraphs suggest a
natural way to construct the space of states $\Ham$ and an algebra of operators, in
terms of the Euclidean space correlation functions, as follows.

Let us again split the whole space $\R^d$, with the coordinates
$x^\mu=(\xx,\tau)$, into two half-spaces
\[
\R^d=\R^d_+\cup\R^d_-
\]
separated by the ``equal time'' slice, the hypersurface $\tau=0$. Namely
\[
\R^d_+:\ (\xx,\tau)\,,\ \tau\ge0\,,\qquad\R^d_-:\ (\xx,\tau)\,,\ \tau\le0\,,
\]
We associate states with the product of field insertions
$\varphi(x_1)\cdots\varphi(x_M)$ at the points $x_1,x_2,\dots,x_M\in\R^d_-$,
\begin{equation}\label{eq:states}
\varphi(x_1)\cdots\varphi(x_M)\ \to\ \ket{\varphi(x_1)\cdots\varphi(x_M)}\,.
\end{equation}
At this point, the above equation just provides notation for the states. We define
$\Ham$ as the linear envelope of the above states
\[
\Ham=\operatorname{Span}\bigl\{\ket{\varphi(x_1)\cdots\varphi(x_M)}:\ x_i\in\R^d_-\bigr\}
\]
The intuition behind this construction comes from the simple minded idea of $\Ham$
as the space of functionals
\[
\Ham=\text{Space of }\Psi[\varphi(\xx)]
\]
of the field configurations $\varphi(\xx)$, $\xx\in\R^{d-1}$. Let
$x_i=(\xx_i,\tau_i)\in\R^d_-$, and consider the functional integral
\begin{equation}\label{eq:Psiint}
\int_{\varphi(\xx,0)=\varphi(\xx)}[D\varphi(x)]\,\varphi(x_1)\cdots\varphi(x_M)\,e^{-\mathcal A[\varphi(x)]}
\end{equation}
over all field configurations $\varphi(x)$ in the half-space $x\in\R^d_-$. This
integral defines (up to overall normalization) the functional
\[
\Psi_{\varphi(x_1)\cdots\varphi(x_M)}[\varphi(\xx)]
\]
which may be regarded as the Schroedinger wave function representing the state
\eqref{eq:states}. In particular, the integral \eqref{eq:Psiint} with no insertions
represents the vacuum state,
\[
\Psi_0[\varphi(\xx)]=\int_{\varphi(\xx,0)=\varphi(\xx)}[D\varphi(x)]\,e^{-\mathcal A[\varphi(x)]}\ \to\ \ket0\in\Ham\,.
\]
The problem with this representation is in the definition of the scalar products.
The simple minded definition
\begin{equation}\label{eq:naiveprod}
\braket{\Psi_2}{\Psi_1}:=\int[D\varphi(\xx)]\,\Psi_2^*[\varphi(\xx)]\,\Psi_1[\varphi(\xx)]
\end{equation}
has all required formal properties, but involves the functional integration which,
as we already know, requires regularization and renormalization. In this context
the problem of renormalization looks much harder than that in the functional
integral \eqref{eq:corrPhi} because the functionals $\Psi[\varphi(\xx)]$ do not have
(quasi) local structure, and generally are far more complicated than
$e^{-\mathcal A}$.

Instead, one defines the scalar product of the states \eqref{eq:states} as
\begin{equation}\label{eq:OSprod}
\bigl\langle\varphi(y_1)\cdots\varphi(y_N)\big|\varphi(x_1)\cdots\varphi(x_M)\bigr\rangle
:=\bigl\langle\varphi(\bar y_1)\cdots\varphi(\bar y_N)\,\varphi(x_1)\cdots\varphi(x_M)\bigr\rangle
\end{equation}
Here the ``mirror image'' points $\bar x$ are defined as follows:
\begin{equation}\label{eq:reflection}
\text{``Reflection'':}\qquad x=(\xx,\tau)\ \leftrightarrow\ \bar x=(\xx,-\tau)
\end{equation}
Then, if all the points $x_1,\dots,x_M,y_1,\dots,y_N$ in \eqref{eq:OSprod} are in
$\R^d_-$, the ``mirror'' points $\bar y_1,\dots,\bar y_N$ are all in $\R^d_+$. The
configuration of the points in the correlation function in the r.h.s.\ of
\eqref{eq:OSprod} is exactly as in \eqref{eq:tauorder}, see figure~\ref{fig:slice}
above.

The most important property of the scalar product is its \emph{positivity},
\[
\braket\Psi\Psi>0
\]
for any $\ket\Psi\in\Ham$, $\ket\Psi\neq0$. For the formal representation
\eqref{eq:naiveprod} the positivity is obvious: Eq.~\eqref{eq:naiveprod} with
$\Psi_1=\Psi_2$ is explicitly positive. For the scalar product defined in terms of
the correlation functions this property is postulated as the \textbf{reflection
positivity}
\[
\sum_{N,M}\int_{x_i,y_k\in\R^d_-}\bigl\langle\varphi(\bar y_1)\cdots\varphi(\bar y_N)\,\varphi(x_1)\cdots\varphi(x_M)\bigr\rangle\,f_N^*(y_1,\dots,y_N)\,f_M(x_1,\dots,x_M)\ge0
\]
for any sequence of functions $f_N(x_1,\dots,x_N)$, which expresses positivity of the
norm of the state
\[
\ket\Psi=\sum_N\int_{x_i\in\R^d_-}f_N(x_1,\dots,x_N)\,\ket{\varphi(x_1)\cdots\varphi(x_N)}\,.
\]
More generally, we can associate states with general products of local fields
$O_a(x)$,
\begin{equation}\label{eq:Ostates}
O_{a_1}(x_1)\cdots O_{a_N}(x_N)\ \to\ \ket{O_{a_1}(x_1)\cdots O_{a_N}(x_N)}\in\Ham\,,\qquad x_1,\dots,x_N\in\R^d_-
\end{equation}
according to the same idea
\[
\Psi_{O_1(x_1)\cdots O_N(x_N)}[\varphi(\xx)]=\int_{\varphi(\xx,0)=\varphi(\xx)}[D\varphi(x)]\,O_1(x_1)\cdots O_N(x_N)\,e^{-\mathcal A[\varphi(x)]}
\]
with the scalar product related to their correlation functions similarly,
\[
\bigl\langle O_{b_1}(y_1)\cdots O_{b_M}(y_M)\big|O_{a_1}(x_1)\cdots O_{a_N}(x_N)\bigr\rangle
:=\bigl\langle\bar O_{b_1}(\bar y_1)\cdots\bar O_{b_M}(\bar y_M)\,O_{a_1}(x_1)\cdots O_{a_N}(x_N)\bigr\rangle
\]
where again $x_1,\dots,x_N;y_1,\dots,y_M\in\R^d_-$, and the points
$\bar y_i\in\R^d_+$ are the ``mirror images'' of $y_i$, as in
Eq.~\eqref{eq:reflection}. However, when the fields $O_i(x_i)$ are not all scalars,
the fields $\bar O_b$ are related to $O_b$ by certain ``reflection''. To get an idea
assume that the field $O(x)$ is real-valued vector with the components $O^\mu(x)$.
Then, if $\mu=1,2,\dots,d-1$ represent ``spatial'' components,
\[
\bar O^\mu(\bar x)=O^\mu(\bar x)\,,\qquad\bar O^\tau(\bar x)=-O^\tau(\bar x)\,,
\]
i.e.\ the sign of the ``temporal'' component $O^\tau$ gets flipped upon reflection.
For example,
\[
\overline{\partial_\tau\varphi}(x)=-\partial_\tau\varphi(x)
\]
so that
\[
\bigl\langle\partial_\nu\varphi(y)\big|\partial_\mu\varphi(x)\bigr\rangle=\pm\bigl\langle\partial_\nu\varphi(\bar y)\,\partial_\mu\varphi(x)\bigr\rangle
\]
where the sign $+$ applies when both $\mu$ and $\nu$ are ``spatial'' or both are
``temporal'', otherwise the sign is $-$. Similar rule applies to higher rank tensors.
For complex valued fields the ``reflection'' involves complex conjugation. The
correlation functions of composite fields $O_a$ obey then the suitably generalized
reflection positivity conditions, which follow from the positivity of the norms of
the states \eqref{eq:Ostates}.

As an elementary application, let us show that in a scale invariant (i.e.\ RG fixed
point) theory the mass dimensions $D_\alpha$ of all nontrivial fields
$\Phi_\alpha(x)$ are positive (we denote $\Phi_\alpha(x)$ the scalar fields that
diagonalize the matrix of anomalous dimensions, see Section~18). From the scale
invariance
\[
\bigl\langle\Phi_\alpha(x)\Phi_\alpha(y)\bigr\rangle=\frac{C_\alpha}{|x-y|^{2D_\alpha}}
\]
where $|x-y|$ is the Euclidean distance between the points $x$ and $y$. Let us take
$\tau>0$, so that $x=(\xx,-\tau/2)\in\R^d_-$. Set $y=\bar x=(\xx,\tau/2)$. We have
\[
\braket{\Phi_\alpha(x)}{\Phi_\alpha(x)}=\frac{C_\alpha}{(\tau^2)^{D_\alpha}}\,.
\]
The positivity requires $C_\alpha>0$; one can set $C_\alpha=1$ by adjusting
normalization of $\Phi_\alpha$. Consider the norm
\begin{equation}\label{eq:normdPhi}
\Bigl\langle\frac\partial{\partial x^1}\Phi_\alpha(\xx,-\tau/2)\Big|\frac\partial{\partial x^1}\Phi_\alpha(\xx,-\tau/2)\Bigr\rangle
\end{equation}
where $x^1$ is one of the spatial coordinates in $\xx=(x^1,x^2,\dots,x^{d-1})$. This
norm is expressed as the second derivative of the correlation function
\[
\frac{\partial^2}{\partial x^1\,\partial y^1}\bigl\langle\Phi_\alpha(\xx,\tau/2)\,\Phi_\alpha(\yy,-\tau/2)\bigr\rangle
\]
evaluated at $\yy=\xx$. Elementary calculation yields for \eqref{eq:normdPhi}
\[
\frac{2D_\alpha}{(\tau^2)^{D_\alpha+1}}
\]
which must be positive, hence $D_\alpha\ge0$.

The case $D=0$ is instructive. If the field $\Phi$ has mass dimension $D=0$ it
follows that
\[
\Bigl|\frac\partial{\partial x^1}\Phi(x)\Bigr\rangle=0\,.
\]
It is then easy to prove that any correlation function involving the derivative of
$\Phi(x)$ vanishes
\[
\bigl\langle\partial_\mu\Phi(x)\,O_1(x_1)\cdots O_N(x_N)\bigr\rangle=0\,,\qquad\Rightarrow\qquad\partial_\mu\Phi(x)\simeq0\,.
\]
Then one can show (using cluster factorization, see below) that
\[
\bigl\langle\Phi(x)\,O_1(x_1)\cdots O_N(x_N)\bigr\rangle=\bigl\langle O_1(x_1)\cdots O_N(x_N)\bigr\rangle
\]
i.e.\ the only field with zero mass dimension is the identity operator.
